\section*{Bab \ref{ch:b2:crustal-melting} --- Pelelehan Kerak: Migmatit dan Granit}

\begin{solution}{exo:b2:crustal-melting:1}
(a) Per \qty{100}{\gram}: $n(\ce{Al2O3}) = 17.03/101.96 = \qty{0.16703}{mol}$,
$n(\ce{CaO}) = 0.09183$, $n(\ce{Na2O}) = 0.06783$, $n(\ce{K2O}) = 0.03076$.
A/CNK $= 0.16703/0.19042 = 0.877$; A/NK $= 0.16703/0.09859 = 1.694$.
(b) A/CNK $< 1 <$ A/NK: \omterm{def:b2:crustal-melting:asi}{metaluminus}. Kalsium yang tidak dapat ditampung feldspar masuk ke
silikat kalsium: hornblende (atau klinopiroksen, titanit).
\end{solution}

\begin{solution}{exo:b2:crustal-melting:2}
(a) Muskovit dan garnet menampung aluminium di luar feldspar: \omterm{def:b2:crustal-melting:asi}{peraluminus}; granit ini
dapat berupa \omterm{def:b2:crustal-melting:types}{granit tipe-S}. (b) Hornblende dan titanit menampung kalsium di luar feldspar:
\omterm{def:b2:crustal-melting:asi}{metaluminus}. (c) Egirin dan amfibol natrium menampung natrium tanpa aluminium
feldspar: \omterm{def:b2:crustal-melting:asi}{peralkalin}.
\end{solution}

\begin{solution}{exo:b2:crustal-melting:3}
Satu kilometer batuan memberikan beban $2800\times9.81\times1000 = \qty{2.747e7}{Pa}$, sehingga
$P = 0.2747\,z$ (\unit{kbar}, $z$ dalam \unit{km}). Garis solidusnya $T = 680 - 7.5(P -
2) = 695 - 2.060\,z$. Menyamakan $10 + 20z = 695 - 2.060z$ memberikan $z = 685/22.060 =
\qty{31.1}{km}$, $T = \qty{631}{\celsius}$, $P = \qty{8.5}{kbar}$.
\end{solution}

\begin{solution}{exo:b2:crustal-melting:4}
(a) Lapisan putih adalah \omterm{def:b2:crustal-melting:migmatite}{leukosom}, jalur hitam adalah \omterm{def:b2:crustal-melting:migmatite}{melanosom}, gneis yang biasa adalah
mesosom. (b) \omterm{def:b2:crustal-melting:migmatite}{Leukosom}, \omterm{def:b2:crustal-melting:migmatite}{melanosom}, dan mesosom yang terpisah dan berbatas jelas: \omterm{def:b2:crustal-melting:migmatite}{metateksit},
dengan struktur stromatik. (c) Batuan itu telah menjadi \omterm{def:b2:crustal-melting:migmatite}{diateksit} (schlieren dan nebulitik):
pelelehan berlanjut lebih jauh, mineral gelap ikut meleleh, dan lelehannya tetap di tempat,
sehingga menghapus perlapisan.
\end{solution}

\begin{solution}{exo:b2:crustal-melting:5}
(a) $M = 39.10 + 3\times24.31 + 26.98 + 3\times28.09 + 12\times16.00 + 2\times1.008 =
\qty{417.30}{\gram\per\mole}$; $w = 18.02/417.30 = \qty{4.32}{\percent}$.
(b) $F = 0.25\times0.0432/0.040 = 0.270$, atau \qty{27}{\percent} massa lelehan: berada dalam rentang
20 sampai \qty{30}{\percent} volume yang dihasilkan percobaan \omterm{def:b2:crustal-melting:reactions}{pelelehan dehidrasi} biotit pada
batuan serupa.
\end{solution}

\begin{solution}{exo:b2:crustal-melting:6}
(a) K-feldspar \qty{278.35}{\gram\per\mole}, albit 262.24, anortit 278.22, muskovit
398.33. (b) Jumlah zat: K-feldspar $30/278.35 = \qty{0.10778}{mol}$, albit $0.13347$,
anortit $0.01797$, muskovit $0.01255$. K-feldspar adalah $\tfrac12$\ce{K2O.}$\tfrac12$\ce{Al2O3},
albit $\tfrac12$\ce{Na2O.}$\tfrac12$\ce{Al2O3}, anortit \ce{CaO.Al2O3}, muskovit
$\tfrac12$\ce{K2O.}$\tfrac32$\ce{Al2O3} per satuan rumus. $n(\ce{Al2O3}) = 0.05389 +
0.06673 + 0.01797 + 0.01883 = \qty{0.15742}{mol}$; $n(\ce{K2O}) = 0.05389 + 0.00628 =
0.06017$, $n(\ce{Na2O}) = 0.06673$, $n(\ce{CaO}) = 0.01797$, jumlahnya $0.14487$.
A/CNK $= 0.15742/0.14487 = 1.087$. (c) Tanpa muskovit, ketiga feldspar memberikan
tepat 1 (\cref{prop:b2:crustal-melting:feldspar}); muskovit, dengan satu \ce{Al2O3}
lebih banyak daripada K-feldspar per satuan rumus, membuat batuan itu \omterm{def:b2:crustal-melting:asi}{peraluminus}.
\end{solution}

\begin{solution}{exo:b2:crustal-melting:7}
Basal: $900 + 2.5\times(1200 - 850) = \qty{1775}{\mega\joule\per\cubic\metre}$. Kerak:
$2.5\times(850 - 650) + 0.20\times900 = \qty{680}{\mega\joule\per\cubic\metre}$.
$V_c/V_b = 1775/680 = 2.6$. Taksiran ini mengabaikan panas yang bocor melalui konduksi
ke atas dan ke samping serta waktu yang diperlukan panas untuk menyebar; taksiran ini juga mengandaikan bahwa seluruh
panas basal mencapai kerak yang siap meleleh: hasilnya adalah batas atas.
\end{solution}

\begin{solution}{exo:b2:crustal-melting:8}
(a) $n(\ce{Al2O3}) = 14.15/101.96 = 0.13878$, $n(\ce{CaO}) = 0.01753$, $n(\ce{Na2O}) =
0.07809$, $n(\ce{K2O}) = 0.04416$; A/CNK $= 0.13878/0.13978 = 0.993$, A/NK $=
0.13878/0.12225 = 1.135$: \omterm{def:b2:crustal-melting:asi}{metaluminus}, nyaris. (b) Dengan \qty{14.30}{\percent} \ce{Al2O3},
$0.14025/0.13978 = 1.003$; dengan \qty{14.00}{\percent}, $0.13731/0.13978 = 0.982$. Indeks itu
mengangkangi angka 1 di dalam ketidakpastian tersertifikasi: obsidian itu terletak di
batas metaluminus--peraluminus, dan kelasnya tidak kokoh.
\end{solution}

\begin{solution}{exo:b2:crustal-melting:9}
(a) Rata-rata $\bar P = \qty{5.0}{kbar}$, $\bar w = \qty{10.045}{\percent}$;
$\sum(P - \bar P)(w - \bar w) = 17.62$, $\sum(P - \bar P)^2 = 14.0$; kemiringan
$17.62/14.0 = \qty{1.259}{\percent}$ per kilobar, titik potong $10.045 - 1.259\times5.0 =
\qty{3.75}{\percent}$. Kemiringan itu sesuai dengan nilai kutipan 1.27 dalam batas
\qty{1}{\percent}. (b) $3.75 + 1.259\times6 = \qty{11.30}{\percent}$. (c) $10/0.1130 =
\qty{88.5}{\gram}$ lelehan. (d) Garis itu bertumpu pada data antara 3 dan \qty{8}{kbar} pada
\qty{800}{\celsius} saja; tidak ada yang menjamin bahwa kelarutan tetap linear pada tekanan
yang lebih tinggi, tempat lelehan dan fluida menjadi makin mirip.
\end{solution}

\begin{solution}{exo:b2:crustal-melting:10}
(a) Garis itu naik \qty{4}{kbar} sepanjang \qty{50}{\celsius}: $P = 6 + 0.08\,(T - 750)$;
pada \qty{770}{\celsius}, $P = \qty{7.6}{kbar}$. (b) Dari 7.6 sampai \qty{6}{kbar},
suhunya turun di sepanjang garis dari 770 sampai \qty{750}{\celsius}; panas yang dilepaskan,
$2.5\times20 = \qty{50}{\mega\joule\per\cubic\metre}$, melelehkan $50/900 = 0.056$, sekitar
\qty{6}{\percent} volume batuan. (c) Inilah jendela yang diajukan untuk
\omterm{def:b2:crustal-melting:leucogranite}{leukogranit} Himalaya, 6 sampai \qty{8}{kbar} dan 750 sampai \qty{770}{\celsius}: dekompresi
saja menghasilkan lelehan, tetapi hanya beberapa persen setiap kali kurva itu dilintasi.
\end{solution}

\begin{solution}{exo:b2:crustal-melting:11}
$\ell^2/\kappa = (2000)^2/10^{-6} = \qty{4e12}{s} = \qty{0.13}{Ma}$ untuk \omterm{def:b2:magma-emplacement:bodies}{sill};
$(10^4)^2/10^{-6} = \qty{1e14}{s} = \qty{3.2}{Ma}$ untuk lapisan underplating. Keduanya singkat
dibandingkan dengan sepuluh juta tahun: satu peristiwa tidak dapat menjaga kerak tetap meleleh selama itu, yang memerlukan
pasokan panas berulang-ulang, melalui intrusi berulang atau melalui pemanasan lambat kerak yang
menebal.
\end{solution}

\begin{solution}{exo:b2:crustal-melting:12}
(a) $5.0x + 2.0(1 - x) = 2.75$ memberikan $x = 0.75/3.0 = 0.25$: seperempat neosom berupa
\omterm{def:b2:crustal-melting:migmatite}{leukosom}. (b) $0.25\times0.20 + 0.75\times4.0 = 0.05 + 3.00 = \qty{3.05}{\percent}$, sama seperti
pada mesosom: tafsiran sistem tertutup berlaku. (c) Untuk \qty{100}{\gram}
batuan induk, terbentuk \qty{25}{\gram} lelehan dan \qty{75}{\gram} residu; jika $m$ gram
lelehan pergi, $(25 - m)/(100 - m) = 0.10$, sehingga $25 - m = 10 - 0.1m$ dan $m = 15/0.9 =
\qty{16.7}{\gram}$: dua pertiga lelehan meninggalkan singkapan itu.
\end{solution}

\begin{solution}{pb:b2:crustal-melting:1}
\noindent\textbf{Bagian I.}

\textbf{1.} $39.10 + 3\times26.98 + 3\times28.09 + 12\times16.00 + 2\times1.008 =
\qty{398.33}{\gram\per\mole}$.

\textbf{2.} $18.02/398.33 = \qty{4.52}{\percent}$.

\textbf{3.} $200\times0.0452 = \qty{9.04}{\gram}$ air per kilogram sekis.

\textbf{4.} $9.97 + (13.82 - 9.97)\times2/3 = \qty{12.54}{\percent}$. Lelehan terkecil adalah
$9.04/0.1254 = \qty{72.1}{\gram}$ per kilogram: $F \ge \qty{7.21}{\percent}$
(\cref{prop:b2:crustal-melting:water}).

\textbf{5.} $9.04/0.060 = \qty{151}{\gram}$: $F = \qty{15.1}{\percent}$.

\textbf{6.} Misalnya: batuan itu mungkin mengandung lebih sedikit muskovit; sebagian muskovit mungkin
sudah terurai di bawah solidus, bersama kuarsa, menjadi K-feldspar dan
aluminosilikat, dengan airnya keluar sebagai fluida sebelum pelelehan dimulai; lelehan mungkin mengandung lebih banyak
air daripada \qty{6}{\percent}; sebagian air mungkin tetap berada dalam biotit.

\noindent\textbf{Bagian II.}

\textbf{7.} $680 - 7.5\times(7.8 - 2) = \qty{636.5}{\celsius}$; puncak batuan itu
\qty{43.5}{\celsius} di bawahnya dan batuan itu tidak dapat meleleh, bahkan dengan air bebas.

\textbf{8.} $605 + (19.2 - 15.3)/(27 - 15.3)\times65 = \qty{626.7}{\celsius}$. Pada
$t = 0$, dasar kerak, pada \qty{595}{\celsius}, berada \qty{32}{\celsius} di bawahnya; pada
\qty{10}{Ma}, pada \qty{656}{\celsius}, dasar itu berada di atasnya.

\textbf{9.} Pelelehan pada solidus basah memerlukan air bebas, dan batuan metamorf yang terpanaskan
hanya mengandung sangat sedikit air bebas: air bebas itu habis dipakai oleh lelehan pertama, yang lalu berhenti bertambah
sampai mika terurai.

\textbf{10.} Antara 30 dan \qty{40}{Ma}: $30 + (800 - 795)/(845 - 795)\times10 =
\qty{31.0}{Ma}$.

\textbf{11.} Dasar kerak turun di bawah \qty{12.5}{kbar} antara 70 dan \qty{80}{Ma}:
$70 + (13.7 - 12.5)/(13.7 - 12.4)\times10 = \qty{79.2}{Ma}$, pada
$914 - 14\times0.923 = \qty{901}{\celsius}$, di dalam rentang amfibol: amfibolit
di dasar kerak akan meleleh.

\textbf{12.} $90 + (11.0 - 10)/(11.0 - 9.6)\times10 = \qty{97.1}{Ma}$, pada $867 -
53\times0.714 = \qty{829}{\celsius}$. Suhu itu berada di atas rentang muskovit: muskovit
telah terurai. Suhu itu berada dalam rentang 800 sampai \qty{850}{\celsius} penguraian biotit
(diukur pada tekanan yang lebih rendah): biotit, mineral hidrus lain dalam sekis itu, berada di
batasnya.

\noindent\textbf{Bagian III.}

\textbf{13.} $n(\ce{Al2O3}) = 15.0/101.96 = \qty{0.14712}{mol}$, $n(\ce{CaO}) =
0.55/56.08 = 0.00981$, $n(\ce{Na2O}) = 3.4/61.98 = 0.05486$, $n(\ce{K2O}) = 4.9/94.20 =
0.05202$.

\textbf{14.} A/CNK $= 0.14712/0.11669 = 1.261$.

\textbf{15.} A/NK $= 0.14712/0.10688 = 1.377$; A/CNK $> 1$: \omterm{def:b2:crustal-melting:asi}{peraluminus}.

\textbf{16.} \omterm{def:b2:crustal-melting:asi}{Peraluminus}, berasal dari sekis: \omterm{def:b2:crustal-melting:types}{granit tipe-S}. Kelebihan aluminiumnya berada dalam
muskovit dan garnet (atau kordierit).

\textbf{17.} $0.14712 - 0.11669 = \qty{0.03043}{mol}$ \ce{Al2O3}; sebagai muskovit,
$0.03043\times398.33 = \qty{12.1}{\gram}$ dalam \qty{100}{\gram} granit.

\textbf{18.} Obsidian itu (nyaris) \omterm{def:b2:crustal-melting:asi}{metaluminus}, tanpa kelebihan aluminium seperti lelehan
metasedimen: sumbernya mungkin berbeda. Akan tetapi, satu sampel hanya sedikit membuktikan:
setiap rangkaian granitoid memiliki anggota yang sedikit \omterm{def:b2:crustal-melting:asi}{peraluminus}, dan suatu rangkaian dinilai dari seluruh
trennya.

\noindent\textbf{Bagian IV.}

\textbf{19.} $10^{12}\,\unit{\cubic\metre}\times2650 = \qty{2.65e15}{kg}$.

\textbf{20.} Granit adalah lelehan dikurangi airnya: $2.65\times10^{15}/0.94 =
\qty{2.82e15}{kg}$ lelehan.

\textbf{21.} Hanya $0.151\times0.70$ massa sekis yang pergi sebagai lelehan:
$2.82\times10^{15}/(0.151\times0.70) = \qty{2.67e16}{kg}$.

\textbf{22.} $2.67\times10^{16}/2800 = \qty{9.5e12}{\cubic\metre}$, yaitu
\qty{9.5e3}{\cubic\kilo\metre}, hampir sepuluh kali \omterm{def:b2:magma-emplacement:bodies}{pluton} itu.

\textbf{23.} $9.5\times10^3/5 = \qty{1.9e3}{\square\kilo\metre}$, sebuah bujur sangkar bersisi sekitar
\qty{44}{km}.

\textbf{24.} Dengan $F = \qty{7.21}{\percent}$, sumbernya $15.1/7.21 = 2.1$ kali
lebih besar, sekitar \qty{2.0e4}{\cubic\kilo\metre}. Volume kerak sumber yang meleleh untuk membentuk
\omterm{def:b2:magma-emplacement:bodies}{pluton} itu, dengan \omterm{def:b2:mantle-melting:fraction}{fraksi lelehan} pertanyaan~5, adalah
\textbf{\qty[text-series-to-math=true]{9.5e3}{\cubic\kilo\metre}}.
\end{solution}
